\chapter*{第 3 章\quad 流体运动学·孔珑题选}
\addcontentsline{toc}{chapter}{第 3 章\quad 流体运动学·孔珑题选}
\markboth{第 3 章\quad 流体运动学·孔珑题选}{}

\begin{tishi}
本章题目取自孔珑《工程流体力学（第四版）》第 4 章"流体运动学和流体动力学基础"的课后习题（原书题号 4-1 至 4-35），只选取其中属于\textbf{流体运动学}的部分，按赵琴教材（精讲本）第 3 章"流体运动学"归章，对应 \S 3.1 描述流体运动的两种方法（欧拉法、拉格朗日法）、\S 3.2 流体运动的基本概念（流线、迹线、脉线、恒定流与非恒定流、维数、加速度）、\S 3.3 连续性方程、\S 3.4 流体微团运动的分析。

\textbf{筛选口径：}凡涉及压强、能量方程、伯努利方程、动量方程、作用力的题目一律不收（这些题归本册第 4 章"流体动力学基础"）。依此，本章只收 \textbf{4-2、4-4 至 4-12} 共 10 题。

\textbf{未收题说明：}4-1（绕过圆柱体的平面流动，求驻点位置与柱面上的速度分布，主要考点属赵琴第 7 章势流与旋涡）、4-3（两平行平板间黏性流动的速度分布与流量，主要考点属赵琴第 8 章黏性流体动力学基础）虽含运动学内容，但不属本章考点；4-13 及其后各题均为伯努利方程、动量方程、力矩与功率类动力学题，按上述口径全部不收。每题题首标注\textbf{孔珑原书题号}，题后给出重要度星级、依据与难度；原书影印文本中量纲口径不当或数值可疑之处，在章末按题号列出说明。
\end{tishi}

\begin{ti}\sloppy\emergencystretch=2em

\item[\textbf{4-2}] 已知平面流动的速度分布规律为
\[ \bv=-\dfrac{\Gamma}{2\pi}\dfrac{y}{x^2+y^2}\,\vec i+\dfrac{\Gamma}{2\pi}\dfrac{x}{x^2+y^2}\,\vec j \]
式中 $\Gamma$ 为常数。求流线方程并画出若干条流线。
\xstarfull{5}\quad\yiju{E4 2014 年试题计算题第 1 题（已知二维速度场求流线方程）、E4 2015 年试题计算题第 1 题（同型流线方程题）、E4 2015 年试题计算题第 7 题（由流函数求速度分布与流线方程）、E4 2019 年试题计算题（求驻点位置及过驻点的流线方程）——流线类计算题在历年真题中出现 $\ge3$ 次}\quad\nandu{中}

\item[\textbf{4-4}] 已知流场的速度分布为
\[ \bv=xy^2\vec i-\dfrac{1}{3}y^3\vec j+xy\vec k \]
\begin{xiaowen}
\item 问属几维流动？
\item 求 $(x,y,z)=(1,2,3)$ 点的加速度。
\end{xiaowen}
\xstarfull{4}\quad\yiju{E4 2021 年试题单选第 4 题（已知平面流速场求某点流体质点的加速度）、E4 2021 年试题单选第 5 题（由速度分量判断一/二/三元流）、E4 2016 年试题填空题第 1 题（已知速度场求点 $(1,1)$ 的加速度 $a_x$、$a_y$）、E4 2020 年试题填空题第 1 题（加速度的欧拉表示式）——加速度与维数判断反复以选择/填空小计算出现}\quad\nandu{中}

\item[\textbf{4-5}] 已知流场的速度分布为
\[ \bv=x^2y\vec i-3y\vec j+2z^2\vec k \]
\begin{xiaowen}
\item 问属几维流动？
\item 求 $(x,y,z)=(3,1,2)$ 点的加速度。
\end{xiaowen}
\xstarfull{4}\quad\yiju{E4 2016 年试题填空题第 1 题、E4 2020 年试题填空题第 1 题（矢量形式的加速度欧拉表示式）、E4 2021 年试题单选第 4 题（同型加速度计算）——按欧拉法求加速度的分量公式是本章选择与填空的常考小计算}\quad\nandu{中}

\item[\textbf{4-6}] 已知流场的速度分布为
\[ \bv=(4x^3+2y+xy)\vec i+(3x-y^3+z)\vec j \]
\begin{xiaowen}
\item 问属几维流动？
\item 求 $(x,y,z)=(2,2,3)$ 点的加速度。
\end{xiaowen}
\xstarfull{4}\quad\yiju{E4 2021 年试题单选第 5 题（由三个速度分量判流动元数）、E4 2021 年试题单选第 4 题（已知流速场求质点加速度）、E4 2019 年试题填空题第 3 题（由速度分量求加速度分量 $a_x$、$a_y$）}\quad\nandu{中}

\item[\textbf{4-7}] 有一输油管道，在内径为 20 cm 的截面上流速为 2 m/s，求另一内径为 5 cm 截面上的流速以及管道内的质量流量。已知油的相对密度为 0.85。
\xstarfull{4}\quad\yiju{E4 2021 年试题计算题（质量流量 $q_m=15$ kg/s、$d_1=100$ mm、$d_2=75$ mm 的变径排出管路）、E4 2022 年 818 单选第 4 题（求出口平均密度 $\rho_2$ 与质量流量 $Q_m$）、E4 2015 年试题填空题第 3 题（变直径管 $d_1=320$ mm、$d_2=160$ mm 求 $v_2$）、E4 2020 年试题填空题第 6 题（连续性方程表示控制体的什么守恒）}\quad\nandu{易}

\item[\textbf{4-8}] 在一内径为 5 cm 的管道中，流动的空气的质量流量为 0.5 kg/s，在某一截面上压强为 $5\times10^5$ Pa，温度为 100 ℃。求在该截面上的气流平均速度。
\xstarfull{4}\quad\yiju{E4 2016 年试题计算题（已知出口 $p_2=452$ kPa、$A_2=5.12\times10^{-4}$ m$^2$ 与 $R=287$ J/(kg$\cdot$K) 求空气质量流量）、E4 2019 年试题计算题（已知环境压强与 $R=287$ J/(kg$\cdot$K) 求喷管质量流量 $Q_m$）、E4 2022 年 818 单选第 4 题（求出口平均密度 $\rho_2$ 与质量流量 $Q_m$）——气体密度换算与质量流量多次作为计算题的核心步骤}\quad\nandu{中}

\item[\textbf{4-9}] 长 3 cm 的锥形喷嘴，其两端内径分别是 8 cm 和 2 cm，流量为 0.01 m$^3$/s，流体无黏性且不可压缩。试导出沿喷嘴轴向的速度表达式，$x$ 距离从大内径一端的端面计起。
\tuyi{锥形喷嘴是一段直线收缩的圆管：大内径端直径 8 cm（半径 4 cm），小内径端直径 2 cm（半径 1 cm），轴向长 3 cm；半径沿轴向线性减小，半锥角为 $45^\circ$；$x$ 自大内径端面沿轴线量起。}
\xstarfull{3}\quad\yiju{E4 2015 年试题填空题第 3 题（变直径管由连续性方程求 $v_2$）——同类一元连续性应用在真题中出现 1 次；E2（第 1 至 3 章偏概念，掌握基本公式与解释即可）}\quad\nandu{中}

\item[\textbf{4-10}] 已知流场内的速度分布为
\[ \bv=\dfrac{4x}{x^2+y^2}\,\vec i+\dfrac{4y}{x^2+y^2}\,\vec j \]
求证：通过任意一个以原点为圆心的同心圆的流量都相等（$z$ 方向取单位长度）。提示：流场速度用极坐标表示。
\xstarfull{3}\quad\yiju{E4 2015 年试题计算题第 8 题（均匀流叠加点源求驻点与流线，涉及点源流场与极坐标表示）、E4 2019 年试题计算题（叠加流场求驻点及过驻点的流线方程）；本考点本身属赵琴第 7 章势流，真题未独立出现}\quad\nandu{中}

\item[\textbf{4-11}] 由空气预热器经两条管道送往锅炉喷燃器的空气质量流量 $q_m=8000$ kg/h，气温 400 ℃，管道截面尺寸均为 400 mm$\times$600 mm。已知标准状态（0 ℃，101325 Pa）空气的密度 $\rho_0=1.29$ kg/m$^3$，求输气管道中空气的平均流速。
\xstarfull{3}\quad\yiju{E4 2016 年试题计算题（空气流量与状态方程换算）、E4 2019 年试题计算题（同型气体质量流量换算）、E4 2022 年 818 单选第 4 题（求气体断面平均密度）——真题以单管形式出现，本条为双管送风，故按三星处理}\quad\nandu{中}

\item[\textbf{4-12}] 比体积 $v=0.3816$ m$^3$/kg 的汽轮机废汽沿一直径 $d_0=100$ mm 的输气管进入主管，质量流量 $q_m=2000$ kg/h，然后沿主管上的另两支管输送给用户。已知用户的需用流量分别为 $q_{m1}=500$ kg/h、$q_{m2}=1500$ kg/h，管内流速均为 25 m/s。求输气管中蒸汽的平均流速及两支管的直径 $d_1$、$d_2$。
\tuyi{习题 4-12 图（原书图 4-26）：主管自左侧水平进入，其上分出两根支管分别通向两个用户；主管断面直径为 $d_0$、质量流量为 $q_m=2000$ kg/h；两支管的质量流量分别为 $q_{m1}=500$ kg/h 与 $q_{m2}=1500$ kg/h，支管内流速均为 25 m/s，待求直径为 $d_1$、$d_2$。}
\xstarfull{4}\quad\yiju{E4 2021 年试题计算题（由质量流量与管径关系求变径管路的流速与管径）、E4 2019 年试题计算题（由连续性方程 $Q=Q_1+Q_2$ 处理分流）、E4 2022 年 818 单选第 4 题（质量流量 $Q_m$ 与密度换算）}\quad\nandu{中}

\end{ti}

\begin{banfa}
\textbf{按题号给出的应试提示}\\
\textbf{第 4-2 题：}流线用 $\dfrac{\dd x}{v_x}=\dfrac{\dd y}{v_y}$（三维写 $\dfrac{\dd x}{v_x}=\dfrac{\dd y}{v_y}=\dfrac{\dd z}{v_z}$）。
代入后先把两个分量中的\textbf{公共因子约掉}，再分离变量积分；本例得 $x^2+y^2=C$，即同心圆族。
流向取圆上一点代入速度符号判断（本例 $\Gamma>0$ 时逆时针）。\\
\textbf{第 4-4、4-5、4-6 题：}两步走。
① \textbf{判维数}：看三个速度分量里出现了哪几个空间坐标，出现了几个就是几维；
判据是"运动要素是否依赖该坐标"，与"该方向有没有速度分量"无关。
② \textbf{求加速度}：先写出九个一阶偏导，再按
$a_x=\dfrac{\partial v_x}{\partial t}+v_x\dfrac{\partial v_x}{\partial x}+v_y\dfrac{\partial v_x}{\partial y}+v_z\dfrac{\partial v_x}{\partial z}$
逐项代入；题给速度场不显含 $t$，时变项直接删去。分数结果化成小数便于校核。\\
\textbf{第 4-7、4-8、4-11、4-12 题：}一元总流连续性——
不可压缩流体 $v_1A_1=v_2A_2$（即 $v_1d_1^2=v_2d_2^2$，流速与内径平方成反比）；
可压缩流体用 $q_m=\rho_1v_1A_1=\rho_2v_2A_2$，气体密度由 $\rho=\dfrac{p}{RT}$ 或标准状态换算
$\rho_1=\rho_0\dfrac{T_0}{T_1}$ 得到；分流（汇流）管路按"进等于出"写 $\sum q_{mi}=\sum q_{mj}$；
求管径用 $d=\sqrt{\dfrac{4q_m}{\pi\rho v}}$。代入前把 kg/h 化成 kg/s、℃ 化成 K、mm 与 cm 化成 m。\\
\textbf{第 4-9 题：}沿程速度分布三步——① 由几何关系写出 $d(x)$；② 写断面面积 $A(x)=\dfrac{\pi d^2}{4}$；
③ 用 $v(x)=\dfrac{q_V}{A(x)}$。最后把两端点代回校核（面积小的一端速度最大）。\\
\textbf{第 4-10 题：}见到 $\dfrac{x}{x^2+y^2}$、$\dfrac{y}{x^2+y^2}$ 一律换极坐标（$x^2+y^2=r^2$），
算出径向速度 $v_r$ 与切向速度 $v_\theta$；穿过半径 $r$ 的圆柱面（$z$ 向取单位长度）的流量为
$q_V=v_r\cdot2\pi r$，若结果与 $r$ 无关即为点源。
\end{banfa}

\begin{tishi}
\textbf{章末说明（按孔珑原书题号）}\\
\textbf{第 4-9 题：}原书解答把 $d=8-2x$（$x$ 以 cm 计）直接代入 $q_V=0.01$ m$^3$/s，写为
$v=\dfrac{0.01}{\pi(4-x)^2}$，量纲上省略了 $10^{-4}$ 的面积换算；本册按统一单位给出
$v=\dfrac{100}{\pi(4-x)^2}$ m/s（$x$ 以 cm 计），两端点流速为 1.99 m/s 与 31.8 m/s，可用题给流量校核。\\
\textbf{第 4-11 题：}原书结果印为 8.987 m/s；按 $\rho_1=\rho_0\dfrac{T_0}{T_1}=1.29\times\dfrac{273}{673}=0.5233$ kg/m$^3$
复算为 8.85 m/s，相差约 1.6\%，疑为原书算术笔误，本册按公式计算值给出。\\
\textbf{第 4-12 题：}原书取支管内流速为 27 m/s（误用主管流速）代入，得 $d_1=0.05$ m、$d_2=0.0866$ m；
按题给"管内流速均为 25 m/s"应为 $d_1\approx0.052$ m、$d_2\approx0.090$ m，本册按题给数据计算。
\end{tishi}

\begin{yicuo}
\textbf{本章四个高频陷阱：}\\
① \textbf{恒定流与均匀流不是一回事。}恒定看\textbf{时间}（$\dfrac{\partial(\cdot)}{\partial t}=0$），
均匀看\textbf{空间}（流线是相互平行的直线）；恒定流可以是非均匀流，非恒定流也可以是均匀流。\\
② \textbf{判维数看"运动要素依赖哪几个坐标"，不看"某方向有没有速度"。}
第 4-4 题中 $v_z=xy$ 并不为零，但它与 $z$ 无关，仍是二维（平面）流动；
反之第 4-6 题中 $v_y$ 含 $z$，即使该方向加速度为零，也是三维流动。\\
③ \textbf{求加速度别漏项。}每个分量都要依次乘以三个速度分量再乘对应的偏导，
如第 4-4 题的 $a_z$ 必须包含 $v_y\dfrac{\partial v_z}{\partial y}$ 这一项，漏掉就会把
$\dfrac{2}{3}xy^3$ 错算成 $xy^3$。\\
④ \textbf{单位陷阱。}$q_m$ 给的是 kg/h、气温给的是 ℃、管径给的是 mm 或 cm，
代入前一律换成 kg/s、K、m；气体还要先用 $\rho=\dfrac{p}{RT}$ 把密度算出来。
\end{yicuo}
